\documentclass[11pt]{amsart}
\usepackage[T1]{fontenc}
\usepackage{lmodern}
\usepackage{amsmath,amssymb,amsthm,booktabs,xcolor}
\usepackage[colorlinks=true,linkcolor=blue!55!black,citecolor=blue!55!black,urlcolor=blue!55!black]{hyperref}
\newtheorem{theorem}{Theorem}[section]
\newtheorem{lemma}[theorem]{Lemma}
\title[Line-Graph Signature at Order Nine]{Exact Verification of the Line-Graph Signature Inequality at Order Nine}
\author{Zijian Zeng}
\address{Institute of Computer Science and Digital Innovation, UCSI University, Kuala Lumpur, 56000, MALAYSIA}
\email{zijianzeng@foxmail.com; 1002266693@ucsiuniversity.edu.my}
\date{31 August 2026}
\hypersetup{pdftitle={Exact Verification of the Line-Graph Signature Inequality at Order Nine},pdfauthor={Zijian Zeng}}
\begin{document}
\begin{abstract}
For a connected simple graph $G$, let $c(G)=|E(G)|-|V(G)|+1$ and let
$s(L(G))$ denote the positive minus negative adjacency inertia of its line
graph. We certify that $2s(L(G))\leq c(G)+1$ for every graph $G$ of order
nine. The verification exhausts all $261{,}080$ connected isomorphism
classes and uses exact rational congruence, with an independent verification
by integer characteristic polynomials and Sturm root counts. There are six
equality classes, all unicyclic, and the inequality is strict on the other
$261{,}074$ classes. This is a finite result; the universal inequality
remains open.
\end{abstract}
\maketitle

\section{Introduction}
The adjacency signature of a graph $H$ is
$s(H)=n_+(A(H))-n_-(A(H))$, with zero eigenvalues omitted. A recent
preprint of Paone and Paone proposes the cyclomatic refinement
\begin{equation}\label{eq:conj}2s(L(G))\leq c(G)+1\end{equation}
for finite simple connected graphs~\cite{paone_2026_21706797}. Their exact
computations reached order eight, while their order-nine traversal was only a
floating-point screening. The earlier constant bound $s(L(G))\leq1$ of
Akbari et al.~\cite{Akbari_2026} is false and the signature of connected line
graphs is unbounded~\cite{francis2026signatureconnectedlinegraphs}; those results do not contradict
\eqref{eq:conj}, whose right side grows with the cycle rank. Related older
signature bounds for line graphs are due to Wang and Fan~\cite{Wang_2014}.

Our result closes precisely the finite order-nine certification gap. No
claim is made here about the universal conjecture or absolute priority beyond
the literature search recorded with the supporting data.

\section{Exact reduction}
All graphs below are finite, undirected and simple. Let $G$ be connected,
with $n$ vertices and $m$ edges, and let $B$ be its unsigned vertex--edge
incidence matrix. Put $Q=D+A=BB^{\mathsf T}$.

\begin{lemma}\label{lem:transfer}
For every connected $G$ with at least one edge,
\[s(L(G))=\operatorname{sig}(Q(G)-2I)-c(G)+1.\]
\end{lemma}
\begin{proof}
Entrywise multiplication gives
$B^{\mathsf T}B=A(L(G))+2I$. The nonzero spectra of
$B^{\mathsf T}B$ and $BB^{\mathsf T}=Q$ agree, including multiplicity.
If $z_Q$ is the nullity of $Q$, the additional zero eigenvalues of
$B^{\mathsf T}B$ have multiplicity $m-n+z_Q$ and become eigenvalues $-2$ of
$A(L(G))$. Counting eigenvalues of $Q$ above, at, and below $2$ gives the
formula; the $z_Q$ terms cancel. This also covers the bipartite case, where
$z_Q=1$.
\end{proof}

If $a$ and $z$ count the eigenvalues of $Q$ above and equal to $2$, the same
bookkeeping yields
\begin{equation}\label{eq:threshold}
s(L(G))=2a+z-m,\qquad 4a+2z\leq3m-n+2.
\end{equation}
The crude estimate $2a+z\leq2n$ proves \eqref{eq:conj} whenever
$3m\geq5n-2$. Thus at order nine the range $m\geq15$ also has a short
noncomputational proof. The census nevertheless checks every edge range.

\section{Certified finite theorem}
\begin{theorem}\label{thm:main}
If $G$ is a connected simple graph with $|V(G)|=9$, then
$2s(L(G))\leq c(G)+1$. Exactly six isomorphism classes attain equality.
\end{theorem}
\begin{proof}
Using \texttt{nauty-geng -cd1D8 9 8:36} from the nauty~2.7r3 distribution
\cite{McKay_2014,NautyGuide}, we generated one graph6 representative from
each isomorphism class of connected simple graphs on nine vertices. The
resulting $261{,}080$ records agree with McKay's published online census
\cite{McKayData}. For each representative, the first verifier formed the integer
symmetric matrix $Q-2I$ and computed its inertia by congruence over
$\mathbb Q$. A nonzero diagonal pivot was removed by its exact Schur
complement. If every remaining diagonal entry was zero but an off-diagonal
entry was nonzero, the corresponding two-dimensional principal block was a
hyperbolic plane and contributed one positive and one negative square. An
all-zero remainder contributed its dimension to the nullity. These steps are
congruences, terminate by decreasing dimension, and give the exact inertia,
including singular cases. Lemma~\ref{lem:transfer} then evaluated the integer
slack $\delta(G)=c(G)+1-2s(L(G))$. The result was $\delta=0$ for six records
and $\delta>0$ for the other $261{,}074$ records.

Independently, a second verifier reparsed the graph6 data, recomputed the
integer characteristic polynomial of every $Q-2I$, square-free factored it,
and used a rational Sturm sequence on each factor, restoring its
multiplicity. It reproduced the input digest, graph count, and outcome
counts. Thus every isomorphism class is covered and every sign decision is
exact.
\end{proof}

Table~\ref{tab:edgecounts} gives an audit-friendly aggregation. The entries
sum to $261{,}080$; equality occurs only at $m=9$.
\begin{table}[ht]
\centering\scriptsize
\caption{Number of connected order-nine classes by edge count.}\label{tab:edgecounts}
\begin{tabular}{crrrrrrrrrr}\toprule
$m$&8&9&10&11&12&13&14&15&16&17\\\midrule
count&47&240&797&2075&4495&8404&13855&20303&26631&31400\\\bottomrule
\end{tabular}
\vspace{4pt}
\begin{tabular}{crrrrrrrrr}\toprule
$m$&18&19&20&21&22&23&24&25&26\\\midrule
count&33366&31996&27764&21817&15558&10096&5984&3247&1635\\\bottomrule
\end{tabular}
\vspace{4pt}
\begin{tabular}{crrrrrrrrrr}\toprule
$m$&27&28&29&30&31&32&33&34&35&36\\\midrule
count&770&344&148&63&25&11&5&2&1&1\\\bottomrule
\end{tabular}
\end{table}

\section{Equality certificates and reproducibility}
The six equality representatives are the following literal graph6 strings:
\begin{verbatim}
H?`@C`w  H?`@E_w  H?`DA`g  H?`DA_w  H?`D@`g  H?bB@_W
\end{verbatim}
Each has $(n,m,c)=(9,9,1)$, inertia
$\operatorname{In}(Q-2I)=(5,0,4)$, and hence $s(L(G))=1$. The certificate
file records the six integer characteristic polynomials. A third check
constructs each line graph explicitly and applies exact characteristic-
polynomial root isolation directly to $A(L(G))$.

The archival package includes the literal graph6 census, generator and timing
logs, full verifiers, the direct equality checker, exact outputs, and a file
manifest. The complete command sequence is in \texttt{REPRODUCE.md}.
Computation used no random choices and no floating-point eigenvalue
classification.

\section{Limitations}
Theorem~\ref{thm:main} is a certified finite result, not a proof of
\eqref{eq:conj} for arbitrary order. It does not classify equality graphs
structurally beyond order nine. The direct source of the cyclomatic
conjecture is a non-peer-reviewed preprint as of the search date. The author
has reviewed the statement, exact certificates, citations, and final
manuscript and accepts responsibility for them.

\section*{Data availability}
The code, data, exact certificates, logs, and release manifest accompany this
article.
\enlargethispage{5\baselineskip}
\footnotesize
\bibliographystyle{amsplain}\bibliography{references}
\end{document}
